\begin{Definition}
Let $R$ be a discrete valuation ring of mixed characteristic $(0,p)$, i.e., $R$ has chracteristic~0 and the residue field $R/\mathfrak{m}_R$ of $R$ has characteristic~$p$. Let~$S$ be the ring of integers of a~finite extension of the fraction field $\operatorname{Frac}(R)$ of $R$. Then $\mathfrak{m}^e_S=\mathfrak{m}_RS$ for some positive integer~$e$, called the ramification index. When $e=1$, $S$ is called an unramified extension of~$R$. On the other hand, when $e\neq 1$, $S$ is called a ramified extension of~$R$.
\end{Definition}

\begin{Definition}
Let $R$ be a commutative ring with prime characteristic~$p$. The Frobenius endomorphism $\operatorname{Frob}_p\colon R\rightarrow R$ is defined by $\operatorname{Frob}_p(r)=r^p$ for all $r\in R$.
\end{Definition}

If $R$ is a finite field, the Frobenius endomorphism is an automorphism.

\begin{Definition}
Let $A$ be a finite abelian group, $A^*$ be the homomorphisms $\operatorname{Hom}(A,\mathbb{C}^\times)$ from~$A$ to~$\mathbb{C}^\times$ and $f\colon A \rightarrow \mathbb{C}$ be a function. The discrete Fourier transform $\hat{f}\colon A^*\rightarrow \mathbb{C}$ is defined by $\hat{f}(\chi)=\sum_{a \in A}f(a)\chi (a)$ for any~$\chi \in A^*$.
\end{Definition}

There is an inverse discrete Fourier transform, it is described by the following
\[
f(a)=\frac{1}{|A|}\sum_{\chi \in A^*}\hat{f}(\chi)\overline{\chi}(a).
\]
